\documentclass[10pt,twocolumn]{article}
\usepackage[a4paper, left=1.5cm, right=1.5cm, top=2cm, bottom=3cm]{geometry}
\usepackage[T1]{fontenc}
\usepackage[utf8]{inputenc}
\usepackage[italian]{babel}
\usepackage{amsmath}
\usepackage{titling}
\usepackage{amssymb}
\usepackage{caption}
\usepackage{graphicx}
\usepackage{float}
\usepackage{relsize}
\usepackage{amsmath}
\usepackage{titlesec}
\usepackage{ragged2e}
\usepackage{circuitikz}
\usepackage{booktabs}
\usepackage{longtable}
\usepackage{makecell}
\usepackage{pdflscape}
\usepackage{enumitem}
\usepackage{tikz}
\usepackage{physics}
\usepackage{xcolor}
\usepackage[most]{tcolorbox}
\usepackage{tikz-3dplot}
\usepackage{tikz}
\usepackage{ragged2e}
\usepackage{siunitx}
% \usepackage{booktabs}
\usepackage[colorlinks=true, linkcolor=black]{hyperref}  %per rendere l'indice genrale "interattivo"
\usepackage{enumitem}  %distanza degli itemize
\setlist[itemize]{itemsep=4pt, parsep=1pt}
\newtcolorbox{nota}{
  blanker,
  before skip=1em,
  after skip=1em,
  left=1em,
  borderline west={1pt}{0pt}{black},
  fontupper=\itshape,
  before upper={\noindent\textbf{Nota}:\quad}
}

\setlength{\intextsep}{8pt} % per regolare la distanza grafico tabelle
\setlength{\columnsep}{60pt}
\setlength{\tabcolsep}{4pt}

\begin{document}
\justifying
	\title{\textbf{Misura della curva volt-amperometrica di una lampadina a filo di tungsteno}}
	\author{Tedeschini Michele \hspace{0.7com} Mirolo Manuele \hspace{0.7cm} Stroili Emanuele}
	\date{16 Ottobre 2026}
	\maketitle
	\newgeometry{left=3cm, right=3cm, top=4cm, bottom=4cm}
	\onecolumn
	\tableofcontents
\vspace{3cm}
	\begin{abstract}
		\centering
		\large
    L'esperimento consiste nel verificare la legge delle lenti sottili, graficando la distanza tra lo schermo $q$, e la lente con la distanza tra la lente e la sorgenete $p$. 
       
	\end{abstract}

	\newpage
\restoregeometry
\twocolumn



\section{Legge fisica da verificare}
\section{Apparato sperimentale}
\begin{figure}[H]
\centering
\begin{tikzpicture}[scale=1.2, transform shape, >=stealth]

% ------------------------------------------------
% Banco ottico (binario)
% ------------------------------------------------
\draw[thick, gray!80] (0,0) -- (7,0);
\draw[thick, gray!80] (0,-0.1) -- (7,-0.1); % Doppio tratto per dare l'idea del binario

% ------------------------------------------------
% Componenti (Lampadina, Lente, Schermo)
% ------------------------------------------------
% Lampadina
\draw[thick, blue!65!black] (0.5, 0) -- (0.5, 1.5);
\node[above] at (0.5, 1.5) {\footnotesize Lampadina};

% Lente
\draw[thick, blue!65!black] (3.5, 0) -- (3.5, 1.5);
\node[above] at (3.5, 1.5) {\footnotesize Lente};

% Schermo
\draw[thick, blue!65!black] (6.5, 0) -- (6.5, 1.5);
\node[above] at (6.5, 1.5) {\footnotesize Schermo};

% ------------------------------------------------
% Quote (distanze p e q)
% ------------------------------------------------
% Linee di richiamo verticali
\draw[gray!70] (0.5, -0.2) -- (0.5, -0.6);
\draw[gray!70] (3.5, -0.2) -- (3.5, -0.6);
\draw[gray!70] (6.5, -0.2) -- (6.5, -0.6);

% Quota p
\draw[<->, thick] (0.5, -0.4) -- (3.5, -0.4) node[midway, below] {$p$};
% Quota q
\draw[<->, thick] (3.5, -0.4) -- (6.5, -0.4) node[midway, below] {$q$};

\end{tikzpicture}
\section{Analisi dati}
\section{Conclusione}
\section{Dati}

\end{document}